\section{One pipeline, seven constraint types: details}\label{sec:pipeline}
All experiments use the same code: the feasible set is turned into its automaton, then into $A_q$ (numerically and as a Qiskit circuit), then into amplitude amplification and QAOA. The automaton is built here from the truth table ($n\le16$), which is exponential; what is general is its existence and width, and for larger instances it is written from the problem structure, as the counters of \cref{sec:portfolio}. Instances: cardinality with a quadratic risk objective, knapsack with \emph{different weights}, weighted independent set, set packing, exact cover, 3-colouring and quadratic assignment, with five random instances per family. The summary is in \cref{sec:validation}; this appendix has the tables.


\begin{table*}[t]\centering\small\setlength{\tabcolsep}{3.5pt}
\caption{\orig{exact ideal numerics} Amplitude amplification on seven constraint types with the same pipeline (exact numerics, $n=10$ decision variables unless noted; median over the instances, optimal solutions marked). $w$: automaton width; $a_{\rm unif}$, $a_{1/2}$, $a_{q^*}$: probability of the marked set in the uniform state, in $A_{1/2}|0\rangle$ and in the best-biased $A_q|0\rangle$; $k$: optimal number of amplification iterations. Standard Grover must also verify feasibility in its oracle; $A_q$ does not.}\label{tab:gen_grover}
\begin{tabular}{llrrrrrrrrr}\toprule
family & constraint & inst. & $w$ & $|F|/2^n$ & $a_{\rm unif}$ & $a_{1/2}$ & $a_{q^*}$ & $k_{\rm std}$ & $k_{1/2}$ & $k_{q^*}$\\\midrule
cardinality ($\sum x=4$) & count & 5 & 5 & 0.205 & $9.8\!\cdot\!10^{-4}$ & $2.0\!\cdot\!10^{-3}$ & $3.2\!\cdot\!10^{-3}$ & 25 & 17 & 13\\
knapsack (weights $w_i$) & linear ineq. & 5 & 14 & 0.385 & $9.8\!\cdot\!10^{-4}$ & $2.9\!\cdot\!10^{-3}$ & $3.8\!\cdot\!10^{-3}$ & 25 & 14 & 12\\
independent set & pairwise & 5 & 11 & 0.075 & $9.8\!\cdot\!10^{-4}$ & 0.016 & 0.022 & 25 & 6 & 5\\
set packing & overlap $\le1$ & 5 & 6 & 0.024 & $9.8\!\cdot\!10^{-4}$ & 0.031 & 0.035 & 25 & 4 & 4\\
exact cover & overlap $=1$ & 5 & 3 & $2.9\!\cdot\!10^{-3}$ & $9.8\!\cdot\!10^{-4}$ & 0.250 & 0.902 & 25 & 1 & 0\\
3-colouring & one-hot + edges & 4 & 11 & 0.029 & $1.5\!\cdot\!10^{-3}$ & 0.039 & 0.069 & 21 & 3 & 2\\
assignment (permutations) & row/col sums $=1$ & 5 & 5 & 0.012 & $7.8\!\cdot\!10^{-3}$ & 0.625 & 0.950 & 8 & 0 & 0\\
assignment $4\times4$ ($n{=}16$) & row/col sums $=1$ & 1 & 9 & $3.7\!\cdot\!10^{-4}$ & $1.5\!\cdot\!10^{-5}$ & 0.031 & 0.082 & 201 & 4 & 2\\
\bottomrule\end{tabular}\end{table*}

\begin{table*}[t]\centering\small\setlength{\tabcolsep}{3.5pt}
\caption{\orig{exact ideal numerics} QAOA at depth $p=3$ on the same families (exact numerics, median over instances). Penalty: $|+\rangle^{\otimes n}$, X mixer, best of four penalty weights. $A_q$+GM: negation-forced initial state, objective-only phase, Grover mixer (valid for any constraint). $A_q$+NB: same state, mixer over single flips and swaps inside $F$ (numerics only). Quality: normalised objective with infeasible samples counted as 0; $p{=}0$ is $A_q$ alone.}\label{tab:gen_qaoa}
\begin{tabular}{lrrrrrrrrrr}\toprule
& & \multicolumn{3}{c}{P(feasible)} & \multicolumn{4}{c}{quality} & \multicolumn{2}{c}{P(optimum)}\\\cmidrule(lr){3-5}\cmidrule(lr){6-9}\cmidrule(lr){10-11}
family & inst. & penalty & $A_q$+GM & $A_q$+NB & $A_q$ alone & penalty & $A_q$+GM & $A_q$+NB & penalty & $A_q$+GM\\\midrule
cardinality ($\sum x=4$) & 5 & 0.95 & 1.00 & 1.00 & 0.53 & 0.52 & 0.82 & 0.90 & 0.007 & 0.022\\
knapsack (weights $w_i$) & 5 & 0.89 & 1.00 & 1.00 & 0.53 & 0.42 & 0.80 & 0.63 & 0.002 & 0.037\\
independent set & 5 & 0.89 & 1.00 & 1.00 & 0.58 & 0.57 & 0.86 & 0.65 & 0.035 & 0.161\\
set packing & 5 & 0.82 & 1.00 & 1.00 & 0.55 & 0.50 & 0.87 & 0.74 & 0.063 & 0.472\\
exact cover & 5 & 0.21 & 1.00 & 1.00 & 0.43 & 0.13 & 1.00 & 0.50 & 0.069 & 1.000\\
3-colouring & 4 & 0.73 & 1.00 & 1.00 & 0.59 & 0.46 & 0.91 & 0.81 & 0.048 & 0.345\\
assignment (permutations) & 5 & 0.38 & 1.00 & 1.00 & 0.63 & 0.28 & 1.00 & 0.63 & 0.262 & 1.000\\
\bottomrule\end{tabular}\end{table*}

\begin{table*}[t]\centering\small
\caption{\orig{compiled counts and ideal (noiseless) simulation} Generic compiler (automaton $\to$ circuit), $n=6$ decision variables, Qiskit Aer. $q$: total qubits (decisions + automaton registers); CZ and depth of $A_q$ after transpilation to $\{$CZ, $R_z$, $\sqrt X$, $X$, $R_y\}$ without any hand optimisation; errors are the largest deviation of the simulated circuit probabilities from the exact formula over all Grover iterations / the $p{=}2$ Grover-mixer QAOA circuit.}\label{tab:gen_circ}
\begin{tabular}{lrrrrrr}\toprule
family & $q$ & $w$ & CZ($A_q$) & $k_{\rm std}\to k_{A_q}$ & err.\ Grover & err.\ QAOA\\\midrule
cardinality & 14 & 4 & 200 & $6\to3$ & 3e-15 & 1e-16\\
knapsack & 15 & 5 & 297 & $6\to6$ & 8e-15 & 3e-16\\
independent set & 13 & 3 & 107 & $6\to4$ & 4e-15 & 5e-16\\
set packing & 13 & 3 & 97 & $6\to3$ & 2e-15 & 2e-16\\
exact cover & 9 & 2 & 24 & $6\to0$ & 2e-15 & 1e-16\\
3-colouring & 12 & 3 & 88 & $6\to4$ & 2e-15 & 2e-16\\
\bottomrule\end{tabular}\end{table*}

\begin{table}[t]\centering\small\setlength{\tabcolsep}{4pt}
\caption{\orig{exact automaton construction} Width\,--\,completeness dial on knapsack ($n=12$, six instances, median). The automaton tracks the weight in units of $r$ (a lower bound), so forcing is sound but incomplete: $r=1$ is exact, $r$ above the largest weight keeps no information.}\label{tab:gen_relax}
\begin{tabular}{rrrrr}\toprule
$r$ & width & P(feas.) & $a_{\rm opt}$ & $k_{A_q}$ (std.\ 50)\\\midrule
1 & 33 & 1.00 & $9.8\!\cdot\!10^{-4}$ & 25\\
2 & 18 & 0.71 & $7.3\!\cdot\!10^{-4}$ & 30\\
3 & 12 & 0.47 & $4.9\!\cdot\!10^{-4}$ & 35\\
4 & 9 & 0.43 & $3.7\!\cdot\!10^{-4}$ & 42\\
6 & 6 & 0.39 & $2.4\!\cdot\!10^{-4}$ & 50\\
12 & 1 & 0.39 & $2.4\!\cdot\!10^{-4}$ & 50\\
\bottomrule\end{tabular}\end{table}

\begin{table*}[t]\centering\small\setlength{\tabcolsep}{3.5pt}
\caption{\orig{ideal (noiseless) matrix-product-state simulation} The generic compiler at scale (Aer matrix-product states, 2000 shots, automata written from the problem structure). Qubits include the automaton registers (interleaved with the variables); the bond dimension of the final state equals the automaton width. Feasible: share of samples accepted by the automaton. Uniform: share of feasible strings among all $2^n$ (exact). Mean: mean sample objective $\pm$ standard error against the exact expectation under $A_q$ (dynamic programme over the automaton). Best / opt.: best of 2000 shots / exact optimum.}\label{tab:gen_scale}
\begin{tabular}{lrrrrrrrrr}\toprule
family & $n$ & qubits & $w$ & bond & feasible & uniform & mean (MPS) & mean (exact) & best / opt.\\\midrule
cardinality (count) & 24 & 85 & 5 & 5 & 1.000 & $6.3\!\cdot\!10^{-4}$ & $21.3\pm0.1$ & 21.3 & 32 / 36\\
cardinality (count) & 60 & 274 & 11 & 11 & 1.000 & $6.5\!\cdot\!10^{-8}$ & $52.3\pm0.1$ & 52.2 & 73 / 96\\
independent set (band) & 24 & 68 & 3 & 3 & 1.000 & $7.5\!\cdot\!10^{-4}$ & $23.8\pm0.1$ & 23.6 & 37 / 37\\
independent set (band) & 60 & 176 & 3 & 3 & 1.000 & $1.0\!\cdot\!10^{-8}$ & $66.3\pm0.2$ & 66.1 & 100 / 109\\
independent set (band) & 80 & 236 & 3 & 3 & 1.000 & $2.1\!\cdot\!10^{-11}$ & $83.3\pm0.2$ & 82.7 & 118 / 130\\
knapsack & 18 & 77 & 17 & 17 & 1.000 & 0.218 & $70.9\pm0.3$ & 70.9 & 111 / 116\\
knapsack & 30 & 145 & 25 & 25 & 1.000 & 0.156 & $118.7\pm0.3$ & 118.9 & 173 / 185\\
\bottomrule\end{tabular}\end{table*}

\begin{table}[t]\centering\small\setlength{\tabcolsep}{4pt}
\caption{\orig{exact automaton construction} Limit: market split ($m=3$ equality constraints, random integers up to 99, one planted solution). The automaton is the sound relaxation that forces only when a bound makes a value impossible (state = vector of partial sums). Width: distinct states; P(feas.): probability that $A_{1/2}|0\rangle$ ends on a solution (dead ends lose mass); medians over three instances.}\label{tab:gen_limits}
\begin{tabular}{rrrrr}\toprule
$n$ & width & P(feas.) & uniform & $k_{\rm std}\to k_{A_q}$\\\midrule
10 & 21 & 0.062 & $9.8\!\cdot\!10^{-4}$ & $25\to3$\\
12 & 76 & $2.0\!\cdot\!10^{-3}$ & $2.4\!\cdot\!10^{-4}$ & $50\to17$\\
14 & 113 & $3.9\!\cdot\!10^{-3}$ & $6.1\!\cdot\!10^{-5}$ & $100\to12$\\
16 & 341 & $4.9\!\cdot\!10^{-4}$ & $1.5\!\cdot\!10^{-5}$ & $201\to35$\\
18 & 5\,253 & $6.1\!\cdot\!10^{-5}$ & $3.8\!\cdot\!10^{-6}$ & $402\to100$\\
20 & 12\,067 & $1.5\!\cdot\!10^{-5}$ & $9.5\!\cdot\!10^{-7}$ & $804\to201$\\
\bottomrule\end{tabular}\end{table}


\paragraph{Grover (\cref{tab:gen_grover}, \cref{fig:general}a).} The preparation $A_q$ cuts the amplification iterations on all seven constraint types, and the gain follows the tightness of the constraint. With loose constraints ($|F|/2^n\approx0.2$--$0.4$: cardinality, knapsack) the saving is $1.5$--$2\times$ ($25\to17$ and $25\to14$ iterations at $q=\tfrac12$, $13$ and $12$ with the best bias). With tight constraints it is large: independent set $25\to6$, set packing $25\to4$, 3-colouring $21\to3$, exact cover $25\to1$, and for $3\times3$ assignments the optimum is found by measuring $A_q|0\rangle$ directly (success $0.63$ at $q=\tfrac12$). For the $4\times4$ permutations ($n=16$, $24$ of $65\,536$ strings feasible) standard Grover needs $201$ iterations and $A_q$ needs $4$ ($2$ with the best bias), which is the regime where the construction pays most. The bias $q^\star$ is a free parameter that is cheap to tune, since $a_q$ is a polynomial in $q$.

\paragraph{QAOA (\cref{tab:gen_qaoa}, \cref{fig:general}b).} With the Grover mixer, $A_q$ gives 100\,\% feasible samples on every instance, at every depth and with arbitrary angles, as \cref{thm:qaoa} states (the smallest feasible mass over all instances, depths and methods is $1-3\cdot10^{-16}$). Penalty QAOA with the best of four penalty weights has $21$--$95$\,\% feasible samples, falling with the tightness of the constraint, and a lower quality than $A_q$ alone on all seven families (the penalty ansatz loses what the initial state already guarantees). At $p=3$ the Grover mixer beats the penalty ansatz in quality on 34 of 34 instances (median $0.80$--$1.00$ against $0.13$--$0.57$), and the neighbour mixer on 33 of 34. For the $4\times4$ assignment ($16$ qubits, $p=2$), penalty QAOA gives $2.6$\,\% feasible samples and quality $0.014$, against $100$\,\% and $0.82$ with $P(\mathrm{optimum})=0.20$ for $A_q$ with the Grover mixer. Two caveats. First, on exact cover and on assignments the feasible set has $3$--$6$ elements, so even a random feasible sample is optimal with probability between $1/5$ and $2/3$ (exact cover: $1$--$2$ optima among $3$--$5$ feasible strings; $3\times3$ assignments: $2$--$4$ among $6$): these instances show the mechanism, not difficulty; the informative ones are cardinality, knapsack, independent set, set packing and colouring. Second, the neighbour mixer is a \emph{local} mixer and needs a connected move graph: on exact cover and on permutations single flips and swaps do not connect the feasible strings, the mixer does not move the state, and its quality is that of $A_q$ alone. Local mixers are cheap but structure-dependent (XY for counts, as in \cref{sec:portfolio}); the Grover mixer is general but needs two copies of $A_q$ per layer.

\paragraph{Circuits (\cref{tab:gen_circ}).} The compiler turns each automaton into a Qiskit circuit with no problem-specific code. On six families with $n=6$ the circuit reproduces the exact formulas to $8\cdot10^{-15}$ for all Grover iterations and to $5\cdot10^{-16}$ for $p=2$ Grover-mixer QAOA, and the garbage registers are deterministic given $x$. The compiled $A_q$ costs $24$--$297$ CZ gates without any optimisation, an upper bound that hand-written counters reduce considerably (the portfolio circuit, \cref{tab:pf_depth}).

\paragraph{Relaxed automata (\cref{tab:gen_relax}).} Knapsack with a weight tracked only up to a lower bound in units of $r$ reduces the width from $33$ to $1$ while P(feasible) of $A_q|0\rangle$ falls from $1.00$ to the uniform value $0.39$ and the iterations rise from $25$ to the $50$ of standard Grover. The dial is monotone and sound: no feasible string ever loses weight, as \cref{prop:relax} requires. This is the situation for generic CNF, where exact look-ahead is intractable and bounded look-ahead is the relaxation.



\paragraph{Scale (\cref{tab:gen_scale}).} With the automaton written from the structure of the problem (partial count, partial weight, last $d$ bits of a band graph) the same compiler runs on Aer's matrix-product-state simulator up to $60$ variables for a cardinality constraint ($274$ qubits), $80$ for a band independent set ($236$ qubits) and $30$ for knapsack ($145$ qubits). Every one of $15$ instances gives $100$\,\% feasible shots, against a uniform share of feasible strings that falls to $6.5\cdot10^{-8}$ and $2.1\cdot10^{-11}$; the bond dimension of the final state equals the automaton width ($3$--$25$) because the registers are interleaved with the variables; the mean sample objective agrees with the exact expectation under $A_q$ (a dynamic programme over the automaton) within two standard errors in $14$ of $15$ instances (largest deviation $2.6$). These runs validate the compiled circuits at scale. The sampled distribution is that of the classical walk $\mu_q$, so they say nothing about quantum advantage.

\paragraph{A limit (\cref{tab:gen_limits}).} Market split ($m$ equality constraints on random integers, one planted solution) is the opposite case. Exact look-ahead would need to know which partial-sum vectors can still reach $b$, that is, to solve the instance (its minimal automaton is a single path). The practical automaton is the bound-only relaxation, whose width grows from $21$ at $n=10$ to $12\,067$ at $n=20$ ($m=3$; $m=2$ and $m=4$ give $14\,240$ and $7\,756$), while P(feasible) of $A_q$ falls to $1.5\cdot10^{-5}$ and the iterations shrink only fourfold ($804\to201$). Each automaton state needs a multi-controlled gate, so the circuit grows faster than the iterations shrink: for such constraints the construction does not pay, which is the width trade-off of \cref{prop:relax} in its harshest form.
